% The questions of the olympiad paper, one \element{section}{...} each.
% Sections: sectionI (objective, +3 / -0.75), sectionII (objective, +5 /
% -1.25, set on each question), numerical (a 3-digit answer, +5).
% Edit them in Questions; add sections in Paper design.

\element{sectionI}{
  \begin{question}{q01}
    The sum of the first 100 positive integers is
    \begin{choiceshoriz}
      \wrongchoice{4950}
      \wrongchoice{5000}
      \correctchoice{5050}
      \wrongchoice{10100}
    \end{choiceshoriz}
  \end{question}
}

\element{sectionI}{
  \begin{question}{q02}
    How many positive divisors does $36$ have?
    \begin{choiceshoriz}
      \wrongchoice{6}
      \wrongchoice{8}
      \correctchoice{9}
      \wrongchoice{12}
    \end{choiceshoriz}
  \end{question}
}

\element{sectionI}{
  \begin{question}{q03}
    If $2^x = 64$, then $x$ equals
    \begin{choiceshoriz}
      \wrongchoice{5}
      \correctchoice{6}
      \wrongchoice{7}
      \wrongchoice{8}
    \end{choiceshoriz}
  \end{question}
}

\element{sectionI}{
  \begin{question}{q04}
    The number of diagonals of a regular octagon is
    \begin{choiceshoriz}
      \wrongchoice{16}
      \correctchoice{20}
      \wrongchoice{24}
      \wrongchoice{28}
    \end{choiceshoriz}
  \end{question}
}

\element{sectionII}{
  \begin{question}{q05}\scoring{e=-1.25,v=0,b=5,m=-1.25}
    Let $f(x)=x^2-5x+6$. The sum of the roots of $f(x)=0$ is
    \begin{choiceshoriz}
      \wrongchoice{$-5$}
      \wrongchoice{$1$}
      \correctchoice{$5$}
      \wrongchoice{$6$}
    \end{choiceshoriz}
  \end{question}
}

\element{sectionII}{
  \begin{question}{q06}\scoring{e=-1.25,v=0,b=5,m=-1.25}
    The remainder when $7^{100}$ is divided by $5$ is
    \begin{choiceshoriz}
      \correctchoice{1}
      \wrongchoice{2}
      \wrongchoice{3}
      \wrongchoice{4}
    \end{choiceshoriz}
  \end{question}
}

\element{numerical}{
  \begin{questionmultx}{nq1}
    Find the smallest positive integer that is divisible by each of $1$, $2$, $3$, $4$, $5$ and $6$.
    \AMCnumericChoices{60}{digits=3,decimals=0,sign=false}
  \end{questionmultx}
}

\element{numerical}{
  \begin{questionmultx}{nq2}
    The sum of the digits of $2^{10}$ is
    \AMCnumericChoices{7}{digits=3,decimals=0,sign=false}
  \end{questionmultx}
}
